% Section 5: Empirical Challenges

\subsection{The Cogitate Collaboration}

The largest empirical test of consciousness theories to date was published in April 2025 by
Ferrante et al.\ in \emph{Nature} \cite{ferrante2025adversarial}. The Cogitate consortium
employed an adversarial collaboration methodology that brought together the principal
architects of Integrated Information Theory (Tononi, Koch) and Global Neuronal Workspace
Theory (Dehaene) as co-authors, alongside prominent consciousness researchers including
Chalmers. This design was intended to maximize scientific rigor by forcing theory proponents
to jointly develop falsifiable predictions \emph{before} data collection
\cite{melloni2023adversarial}.

The study tested 256 participants across three neuroimaging modalities: functional magnetic
resonance imaging (fMRI, $n=120$), magnetoencephalography (MEG, $n=102$), and intracranial
electroencephalography (iEEG, $n=34$). All predictions were preregistered, and the
experimental protocols were designed to pit IIT's predictions against GNWT's in a direct
empirical confrontation.

\subsection{Three Predictions Tested}

The adversarial protocol focused on three key predictions that differentiated IIT from GNWT.
Table~\ref{tab:cogitate} summarizes the theoretical predictions and empirical outcomes.

\begin{table}[htbp]
\centering
\caption{Preregistered predictions and empirical results from the Cogitate adversarial
collaboration \cite{ferrante2025adversarial}.}
\label{tab:cogitate}
\begin{tabular}{p{0.22\linewidth}p{0.34\linewidth}p{0.34\linewidth}}
\toprule
\textbf{Prediction} & \textbf{IIT} & \textbf{GNWT} \\
\midrule
\textbf{Content Location} & Conscious content decodable from posterior cortex; PFC not
necessary & Conscious content requires PFC activation for access \\
\textit{Result} & \multicolumn{2}{p{0.68\linewidth}}{\textit{Partial confirmation for both:
content decodable from posterior and PFC, but PFC encoded only category (not identity/orientation)}} \\
\midrule
\textbf{Temporal Dynamics} & Sustained posterior activity throughout stimulus duration &
Ignition events at onset and offset; silent storage between \\
\textit{Result} & \multicolumn{2}{p{0.68\linewidth}}{\textit{IIT confirmed (sustained
posterior activity). GNWT failed: no ignition at stimulus offset}} \\
\midrule
\textbf{Interareal Connectivity} & Sustained gamma-band synchronization within posterior
cortex & Long-range frontal--sensory connectivity \\
\textit{Result} & \multicolumn{2}{p{0.68\linewidth}}{\textit{IIT failed: only 25/657
posterior electrodes showed sustained synchronization. Some GNWT-consistent
frontal--visual synchrony observed}} \\
\bottomrule
\end{tabular}
\end{table}

\textbf{Prediction 1---Content Location.} IIT predicted that conscious content should be
decodable from posterior cortical regions (the ``posterior hot zone''), with prefrontal
cortex (PFC) playing no necessary role \cite{tononi2016integrated}. GNWT predicted that PFC
activation is necessary for conscious access \cite{dehaene2014toward}. The empirical result
was nuanced: content was decodable from both regions, but PFC represented only stimulus
\emph{category} (e.g., face vs.\ object), whereas posterior cortex encoded fine-grained
details. This complicates both theories' predictions.

\textbf{Prediction 2---Temporal Dynamics.} IIT predicted sustained posterior cortical
activity throughout conscious perception \cite{tononi2004information}. GNWT predicted
transient ``ignition'' events at both stimulus onset and offset \cite{dehaene2014toward}.
The data confirmed IIT's prediction but decisively falsified GNWT's offset prediction: no
ignition event was observed when the stimulus disappeared from consciousness. GNWT
successfully explains access to consciousness but fails to account for its termination
dynamics.

\textbf{Prediction 3---Interareal Connectivity.} IIT's strongest prediction concerned
sustained gamma-band (30--100 Hz) synchronization within posterior cortex
\cite{tononi2016integrated, oizumi2014phenomenology}. GNWT predicted long-range
frontal--sensory connectivity. The iEEG data provided the highest resolution test. IIT's
prediction was decisively falsified: only 25 of 657 posterior cortical electrodes exhibited
sustained gamma synchronization. Koch, a principal architect of IIT, acknowledged:
``The challenge is to come up with ways to explain the null result in the third prediction''
\cite{ferrante2025adversarial}.

\subsection{Implications for AI Consciousness}

The Cogitate results carry profound implications for the AI consciousness debate. Both IIT
and GNWT have been deployed as adjudicative frameworks in recent assessments
\cite{butlin2023consciousness, shin2025iit, goldstein2024case}. Yet neither theory survived
empirical scrutiny intact.

IIT's failure on the connectivity prediction is particularly damaging to its anti-AI
arguments. Tononi and Koch have long maintained that artificial neural networks cannot possess
high $\Phi$ because they lack the dense causal integration of biological brains
\cite{tononi2016integrated}. But if biological brains themselves do not exhibit the predicted
posterior synchronization---the neural signature of integration---then the metric's validity
is called into question. Are we measuring the right thing?

GNWT's failure on offset dynamics is equally troubling. The theory can explain how
information enters the global workspace (via ignition), but it cannot explain how it exits.
For AI systems, this means that satisfying GNWT criteria (broadcast architecture, attention
mechanisms \cite{vaswani2017attention}) may correlate with certain functional capacities but
does not guarantee phenomenal consciousness.

The meta-implication is stark: \textbf{we are attempting to measure AI consciousness with
instruments that cannot even fully explain biological consciousness.} If our best
neuroscientific theories fail to predict the neural correlates of consciousness in humans,
on what grounds do we deploy those same theories to adjudicate consciousness in machines?

\subsection{The Measurement Crisis}

The Cogitate results illuminate a broader measurement crisis in consciousness science. Three
interrelated problems afflict current approaches:

\textbf{Computational Intractability.} Computing $\Phi$ is NP-hard and becomes intractable
for networks with more than approximately 10 nodes \cite{albantakis2023iit4}. Shin et al.\
attempted to sidestep this limitation by using perplexity as a proxy for $\Phi$
\cite{shin2025iit}. But perplexity is a behavioral measure of prediction accuracy, not a
structural measure of causal integration. This substitution smuggles in a functionalist
assumption---that behavioral performance tracks consciousness---precisely the assumption IIT
was designed to reject.

\textbf{Operationalization Gaps.} GNWT's criteria---global broadcasting, attentional
selection, reportability---lack precise operationalization for artificial systems
\cite{dehaene2014toward}. What constitutes ``broadcasting'' in a transformer architecture
\cite{vaswani2017attention}? Is self-attention a form of global access? These questions have
no agreed-upon answers. Goldstein et al.\ argue that LLMs satisfy GNWT criteria because they
exhibit global information integration via attention \cite{goldstein2024case}. But this
interpretation is contested: critics note that attention weights are determined by training
dynamics, not by a central executive \cite{butlin2023consciousness}.

\textbf{Theory-Ladenness of Evidence.} Every proposed marker of consciousness presupposes
a theory. IIT's $\Phi$ presupposes that integration is the essence of consciousness. GNWT's
broadcast criterion presupposes that access is constitutive of phenomenality. FEP approaches
presuppose that consciousness is variational inference \cite{friston2010free, wiese2024fep}.
There are no theory-neutral empirical markers. This reflects a deeper epistemological
problem identified by Marr \cite{marr1982vision}: different levels of analysis may be
incommensurable.

The path forward requires theory-independent empirical approaches. The adversarial
methodology of Cogitate offers a model: force competing theories to make differential
predictions and test them jointly. For AI systems, this might involve designing experiments
where IIT and GNWT make opposite predictions about system behavior, then measuring which
prediction holds. Only through such empirical adjudication can we move beyond philosophical
speculation toward a science of consciousness that encompasses both biological and artificial
minds.
